\documentclass[11pt]{article}
\usepackage[T1]{fontenc}
\usepackage[a4paper,margin=1in]{geometry}
\usepackage{amsmath,amssymb,amsthm}
\usepackage[hidelinks]{hyperref}
\usepackage{microtype}

\newtheorem*{conjecture}{Conjecture 3.6 (Paulsen--Solazzo)}
\newtheorem*{proposition}{Proposition}

\setlength{\parindent}{0pt}
\setlength{\parskip}{0.6em}
\renewcommand{\thesection}{\arabic{section}.}

\title{Literature-Implied Answer to Conjecture 3.6:\
Semialgebraic Hypographs Have Finitely Many Corners}
\author{Run \texttt{fa\_banach\_001} --- lane 1}
\date{22 July 2026}

\begin{document}
\maketitle

\begin{center}
\fbox{\parbox{0.9\textwidth}{
\textbf{Status: full resolution implied by a classical theorem.}
If the hypograph is semialgebraic, then its upper fiber endpoint $f$ is a
semialgebraic function.  Semialgebraic $C^1$ cell decomposition makes $f$
$C^1$ away from finitely many points.  Thus the conjectured statement is true.
}}
\end{center}

\section{Source conjecture}

Paulsen and Solazzo pose the following as Conjecture~3.6 on arXiv PDF
page~26~\cite{source}.

\begin{conjecture}
Let $f:[0,1]\to\mathbb R$ be continuous and nonnegative, and set
\[
 C_f=\{(x,y):0\le x\le1,\ 0\le y\le f(x)\}.
\]
If $f$ is nondifferentiable at infinitely many points, then $C_f$ is not
semialgebraic.
\end{conjecture}

The source says this was the only obstruction to its proposed
nonsemialgebraic operator-algebra ball.

\section{Known regularity theorem}

The $C^1$ cell decomposition theorem for definable functions says that if
$X\subseteq\mathbb R^n$ and $g:X\to\mathbb R$ are definable, then there is a
finite decomposition of $\mathbb R^n$ into $C^1$ cells partitioning $X$ such
that $g$ restricts to a $C^1$ function on every cell contained in $X$.
See Bhardwaj--van den Dries, Theorem~A.14 (PDF page~35)~\cite{bvd}.
Semialgebraic sets and functions are definable in the real field, so the
theorem applies to them.

In dimension one, a finite cell decomposition consists of finitely many
points and finitely many open intervals.  Consequently, every semialgebraic
function on $[0,1]$ is $C^1$ off a finite set.

\newpage
\section{Deduction}

\begin{proposition}
Conjecture~3.6 is true.
\end{proposition}

\begin{proof}
Suppose toward a contradiction that $C_f$ is semialgebraic.  Its upper-fiber
endpoint set is
\[
 \Gamma=\left\{(x,y)\in C_f:
 \neg\,\exists t\in\mathbb R\ \bigl(t>y\ \wedge\ (x,t)\in C_f\bigr)
 \right\}.
\]
Semialgebraic sets are closed under Boolean operations and coordinate
projections (Tarski--Seidenberg), so $\Gamma$ is semialgebraic.  For every
$x\in[0,1]$, the vertical fiber of $C_f$ is exactly $[0,f(x)]$; hence
\[
 \Gamma=\{(x,f(x)):0\le x\le1\}.
\]
Thus $f$ is a semialgebraic function.

Apply $C^1$ cell decomposition to $f:[0,1]\to\mathbb R$.  There is a finite
partition of $[0,1]$ into point cells and interval cells such that $f$ is
$C^1$ on each interval cell.  Therefore every possible point of
nondifferentiability lies among the finitely many point cells and interval
endpoints.  This contradicts the hypothesis that $f$ is nondifferentiable at
infinitely many points.  Hence $C_f$ is not semialgebraic.
\end{proof}

\section{Scope and provenance}

The proof resolves the source conjecture completely, including the case where
$f$ vanishes on a nontrivial set: the maximal-fiber formula defining
$\Gamma$ remains valid.  By the source's Proposition~3.7, its stated
operator-algebra consequence is therefore unconditional.  This note does not
independently audit the $\mathcal O_3/\mathcal O_4$ dimension inconsistency in
the wording and proof surrounding Conjecture~3.5 and Proposition~3.7.

This is not claimed as a new theorem.  It is a direct identification of
Conjecture~3.6 as a consequence of classical semialgebraic regularity.  A
bounded exact-phrase, author/title, citation, local-corpus, and web search
through 22 July 2026 found no later paper explicitly recording this
particular implication.

\begin{thebibliography}{9}
\bibitem{source}
V. I. Paulsen and J. P. Solazzo,
\emph{Interpolation and Balls in $\mathbb C^k$},
J. Operator Theory \textbf{60} (2008), 379--398;
arXiv:math/0604497.

\bibitem{bvd}
N. Bhardwaj and L. van den Dries,
\emph{On the Pila--Wilkie theorem},
arXiv:2010.14046, Appendix A, Theorem~A.14.
\end{thebibliography}

\end{document}
